\documentclass[11pt]{article}
\usepackage[a4paper,margin=1in]{geometry}
\usepackage{setspace,microtype,graphicx,booktabs,tabularx,array,amsmath,amssymb}
\usepackage[round,authoryear]{natbib}
\usepackage[hidelinks]{hyperref}
\usepackage{caption,float}
\usepackage{enumitem}
\usepackage{xspace}
\setstretch{1.12}
\setlength{\emergencystretch}{3em}
\captionsetup{font=small,labelfont=bf}
\newcommand{\Lfourq}{\ensuremath{L_{4,\mathrm{query}}}\xspace}
\newcommand{\code}[1]{\texttt{#1}}

\title{Decision Relevance of Uncertainty Calibration in Active-Learning Materials Screening}
\author{Renjie Li\textsuperscript{1,*}\\
\small \textsuperscript{1}School of Materials and Energy, Southwest University, Chongqing, China\\
\small *Corresponding author: Renjie Li; email: \href{mailto:3037147645@qq.com}{3037147645@qq.com}}
\date{}

\begin{document}
\maketitle

\begin{abstract}
Uncertainty calibration has scientific decision value only through its interaction with candidate selection and the task objective. We separate calibration quality, acquisition sensitivity, candidate perturbation, and sequential utility in a fixed offline materials-screening benchmark. At a fixed candidate state, positive global rescaling leaves the specified rank acquisition unchanged. In the frozen 40-seed confirmation, equal-weight LOCAL--RAW differences in integrated top-1\% recall were 0.0190 (95\% bootstrap CI 0.0122--0.0259) for rank acquisition and 0.0144 (0.0083--0.0206) for magnitude acquisition, while both \code{matbench\_mp\_gap} splits showed losses. A separate 20-seed second randomized-tree-ensemble analysis remained positive at top 1\% on average, but rank was near zero at top 0.5\% and magnitude favored RAW at top 5\%. Retrospective V2 entrant analysis found lower top-1\% enrichment for rank acquisition on MP-gap chemical OOD, but little difference in the other MP-gap policy--split combinations. A separate sequential sensitivity analysis found that the sign and size of LOCAL--RAW utility also varied with the interpolation exponent. These results distinguish a calibration improvement from a useful decision change: neither a better calibration summary nor a changed candidate set alone establishes benefit for the intended discovery objective.
\end{abstract}

\noindent\textbf{Keywords:} materials discovery; active learning; uncertainty calibration; acquisition function; Matbench; decision relevance

\section{Introduction}
\subsection{Sequential decisions in scientific learning}
Scientific active learning allocates a limited computation or experiment budget through a sequence of model-guided queries. A surrogate is fitted to observations available so far; an acquisition policy combines predicted values and uncertainty to select the next candidates. Materials screening offers a useful setting for studying this process because composition--property data support repeated, controlled offline campaigns, while real closed-loop materials studies illustrate how predictions and uncertainty can guide laboratory measurements \citep{dunn2020matbench,kusne2020cameo,rao2022highentropy}. A recent perspective emphasizes that implicit workflow choices and evaluation objectives can limit comparability across materials active-learning studies \citep{nair2026critical}.

\subsection{Prediction-layer calibration is not a decision outcome}
Regression uncertainty can be assessed with interval coverage, calibration error, or proper scoring rules. Post-hoc calibration methods can improve these prediction-layer properties, but the effect depends on the calibration target, evaluation distribution, and uncertainty representation \citep{gneiting2007proper,kuleshov2018calibrated,ovadia2019uncertainty}. Prior work has already connected calibration to sequential optimization. Deshpande et al. integrate online calibration and conformal prediction into Bayesian optimization and report faster convergence on benchmark functions and hyperparameter tasks \citep{deshpande2024online}. Jiang et al. apply global temperature scaling to bootstrap uncertainty and report improved expected-improvement optimization in eight of nine model--alloy-dataset cases \citep{jiang2026calibration_optimization}. In contrast, Harish reports that conformal calibration improved uncertainty reliability but did not improve acquisition when selection depended only on the uncertainty ranking in a rare-earth extraction study; the preferred strategy also depended on how the discovery target was defined \citep{harish2026objective}. Direct materials active-learning studies likewise show objective-dependent outcomes: post-hoc Gaussian-process uncertainty calibration did not improve MEPE campaigns for machine-learned force fields, although a Student-t process could outperform Gaussian processes and random sampling with sufficiently large training sets \citep{thomasmitchell2023calibration}; a recent arXiv preprint reports that calibration often did not reduce the number of labels needed for out-of-distribution model improvement relative to random or uncalibrated sampling \citep{dale2025activelearning}. Utility-directed conformal prediction makes downstream loss explicit, but addresses prediction-set construction rather than sequential materials acquisition \citep{cortesgomez2025utility}. Together, these studies establish that calibration, acquisition behavior, and utility have already been connected in several settings. Our narrower question is how global positive scaling and local candidate-dependent rescaling pass through rank and magnitude acquisition to candidate selection and top-tail utility across fixed Matbench tasks. Calibration metrics alone do not show whether a particular policy changes its selected candidates or whether those candidates advance the specified objective.

The relevant path is therefore \emph{prediction and calibration} $\nightarrow$ \emph{acquisition representation} $\nightarrow$ \emph{selected candidates and campaign trajectory} $\nightarrow$ \emph{discovery utility}. A transformation may improve a coverage summary while leaving a policy's ordering unchanged; a perturbation may change the ordering yet direct queries away from the target tail. This study extends prior materials evidence by testing how rank and magnitude policies transmit global versus candidate-dependent rescaling into selection and target-tail utility across fixed benchmark tasks and splits. Calibration quality, decision perturbation, and utility remain related but non-interchangeable outcomes.

\subsection{Why the acquisition representation matters}
For a rank-based policy, a common positive scale factor preserves every pairwise uncertainty ordering. A candidate-dependent factor can change relative uncertainty and therefore the acquisition score. A magnitude-based policy behaves differently because the numerical scale of uncertainty affects its trade-off with the predicted mean. These structural differences motivate checking policy sensitivity before interpreting downstream campaign results.

\subsection{Contributions and scope}
Using a prespecified offline benchmark over four Matbench regression tasks, IID-record and chemical-system OOD splits, and 40 paired seeds, we provide a policy-focused complement to prior calibration-driven materials-optimization studies. We measure shared-state calibration, candidate-selection changes, and sequential target-tail utility separately; state the positive-global-scaling invariance for the specified rank rule and verify its implementation; and compare candidate-dependent rescaling with RAW and two prespecified controls. The frozen primary pooled utility is positive, but both MP-gap splits are negative. A second randomized-tree-ensemble analysis and retrospective entrant-tail diagnostics expose policy and objective dependence. The study concerns offline materials screening with a composition-based representation and does not evaluate a laboratory closed loop.

\begin{figure}[tbp]
\centering
\includegraphics[width=0.94\linewidth]{figures/decision_relevance_framework.pdf}
\caption{Decision-relevance framework. RAW, GLOBAL, and LOCAL modify uncertainty before rank or magnitude acquisition; PERMUTED and MATCHED provide prespecified comparison conditions. The framework treats calibration quality, policy sensitivity and candidate changes, and target-tail utility as separate evidence stages. The positive-scale invariance applies to rank acquisition under a fixed candidate state and deterministic ties.}
\label{fig:framework}
\end{figure}

\section{Methods}
\subsection{Tasks and experimental design}
The confirmation covered \code{matbench\_expt\_gap}, \code{matbench\_mp\_gap}, \code{matbench\_log\_gvrh}, and \code{matbench\_log\_kvrh} \citep{dunn2020matbench}, each evaluated using an IID-record split and a chemical-system-disjoint OOD split. This four-task set was inherited from the frozen benchmark protocol; a separate pre-confirmation task-selection rationale is not recoverable from the archived project materials, which limits the scope of cross-task interpretation. The design used 40 confirmation seeds (926500--926539) and a fixed campaign geometry inherited from the frozen Matbench protocol. The surrogate was an 80-member randomized extra-tree bagging ensemble with a minimum leaf size of two \citep{geurts2006extratrees}. The benchmark objective was to retrieve candidates with the largest recorded property values. For band-gap tasks, this is a defined retrieval objective rather than a claim that larger gaps are preferable for every application. Each campaign began from the protocol-defined labelled source set and selected candidates only from the target pool. After each batch, the selected labels were revealed and added to the training set for the next stage. Unqueried target labels and diagnostic labels were excluded from surrogate and calibration fitting, acquisition scoring, and control matching. The campaign query budget was 20\% of the target pool, allocated over four sequential batches. The design and analysis choices were recorded in local protocol and freeze files before confirmation; an external preregistration identifier has not been established for this draft.

The run matrix passed its prespecified completeness checks. For each seed, outcomes were paired within each of the eight task--split environments and averaged over those environments. The inferential unit was one paired seed-level mean (\(n=40\)); task--split environments were fixed conditions, not independent population replicates. The confidence intervals quantify algorithmic variability over paired seeds for these fixed benchmark environments; they do not quantify uncertainty over a population of materials-property tasks. Detailed record counts and provenance are reported in the Supplementary Information.

\subsection{Feature representation and preprocessing}
The archived adapter documents a Pymatgen composition fallback for the historical feature branch; the original V5 source lock and feature cache are unavailable, so exact identity with the historical runtime cannot be independently re-established. In the independently frozen V2 snapshot, formulas were parsed with Pymatgen 2024.10.29 \citep{ong2013pymatgen}. For each task, the element vocabulary was the lexicographically sorted set of elements appearing in that task snapshot. Each formula was represented by its atomic fractions over that vocabulary, followed by total atom count and the number of distinct elements. This yielded 83 columns for \code{matbench\_expt\_gap} and 86 for each other task. The matrix was dense float64; missing entries were set to zero and non-finite values rejected. No scaling, target transformation, feature selection, constant-column removal, or categorical encoding was applied. The vocabulary was constructed from all formulas in the task snapshot before splitting, without using target labels; it is therefore an unsupervised, transductive schema step rather than a train-only fitted transform. All four tasks used the same implementation, with task-specific element vocabularies. The V2 matrix hashes and column counts are recorded in the data manifest.

\subsection{Acquisition policies and uncertainty interventions}
Let \(\mu_i\) and \(\sigma_i\) denote the predicted mean and ensemble uncertainty for candidate \(i\) in the current pool. The rank-based acquisition score was
\begin{equation}
 A_{\mathrm{rank}}(i)=r(\mu_i)+r(\sigma_i),
 \label{eq:rank}
\end{equation}
where \(r\) is the within-pool percentile rank. The magnitude-based score was
\begin{equation}
 A_{\mathrm{mag}}(i)=\mu_i+\sigma_i.
 \label{eq:mag}
\end{equation}
Higher acquisition scores were selected first, with ties resolved by a stable material identifier. Raw uncertainty (RAW) used \(\sigma_i\) without a calibration multiplier.

\paragraph{Proposition (positive-scale invariance).} For a fixed candidate state, let \(A_i=r(\mu_i)+r(\sigma_i)\), where \(r\) is the within-pool percentile rank, and let \(c>0\). Then replacing every \(\sigma_i\) by \(c\sigma_i\) leaves every score and selected candidate unchanged. \emph{Proof.} Multiplication by a positive constant preserves pairwise order and ties, so \(r(c\sigma_i)=r(\sigma_i)\) for every candidate. Therefore \(A_i'=A_i\); deterministic tie-breaking produces the same ordered candidate list. If the subsequent model update is deterministic, induction over query batches gives the same trajectory. \(\square\)

The local correction was learned from out-of-bag (OOB) residual ratios derived from the currently labelled training data, including labels revealed by preceding queries. For observations with sufficient OOB predictions, define \(q_j=|y_j-\mu_j^{\mathrm{OOB}}|/\max(\sigma_j^{\mathrm{OOB}},10^{-8})\). Observations required at least three OOB predictions, and fitting required at least eight valid residual ratios. An ExtraTrees regressor (64 trees; minimum leaf size five) was fitted to \(\log(q_j+10^{-8})\) using the composition-based features. Five-fold out-of-fold predictions of this log scale were used to estimate a median normalization multiplier, targeting the marginal median of the standardized absolute residual under a Gaussian reference. The predicted local scale was clipped to [0.25, 4]. The confirmation correction used a frozen geometric interpolation with exponent \(\alpha=0.5\) between the global scale and the local scale. The uncertainty supplied to LOCAL was \(\max(\sigma_i,10^{-8})s_{\alpha}(x_i)\). This is a marginal correction and does not guarantee conditional calibration.

The PERMUTED condition independently permuted local scale assignments within each seed and stage, retaining the empirical scale distribution but disrupting its association with candidate features. The MATCHED condition was an acquisition intervention rather than a calibration estimator. For magnitude acquisition, its scalar exploration weight matched the root-mean-square uncertainty bonus on the same remaining pool. For rank acquisition, a fixed grid selected the weight whose batch disagreement from RAW most closely matched the LOCAL-versus-RAW disagreement, with ties resolved by proximity to one and then by the lower weight. Consequently, comparisons with MATCHED address these prespecified exploration summaries, not every possible alternative policy.

For rank acquisition, a GLOBAL condition applied a positive scalar to all uncertainties. Its purpose was to test a mathematical invariance, not to estimate an independent utility benefit. Since \(r(s\sigma_i)=r(\sigma_i)\) for \(s>0\), Eq.~\nef{eq:rank} is unchanged. Under fixed candidate states and deterministic tie handling, the selected batches and resulting trajectory must therefore be identical.

\subsection{Endpoints and statistical analysis}
The primary endpoint, \Lfourq, was the normalized query-order area under the recall curve for the target-pool top 1\%. Let \(N\) be the initial target-pool size, \(K=\max(1,\lceil0.01N\nceil)\), and \(R_b\) the fraction of these \(K\) candidates retrieved after \(b\) queries. The implementation used the trapezoidal endpoint
\begin{equation}
 L_{4,\mathrm{query}}=\frac{1}{B}\sum_{b=1}^{B}\frac{R_{b-1}+R_b}{2},\qquad R_0=0,
\end{equation}
where \(B\) is the total query budget. Top-tail ties were resolved by stable material identifiers. Candidates within a batch were ordered by their acquisition score; model updates occurred between batches, not between individual candidates. Thus this endpoint measures discovery efficiency along the stored query ordering, not experimental elapsed time or a batch-checkpoint-only area. The six prespecified contrasts were LOCAL minus RAW, PERMUTED, and MATCHED for each of the rank and magnitude acquisition rules. A value of 0.005 was used as a protocol-defined descriptive practical threshold; the archived plan does not preserve a scientific rationale for this value, and it is not treated as an externally validated minimal important difference.

For each contrast, the difference was computed within each of the 40 paired seed blocks. A percentile bootstrap with 10,000 resamples produced the 95\% confidence interval. A two-sided sign-flip test with 20,000 draws tested the paired mean difference; Holm correction controlled family-wise error across the six primary tests. The bootstrap and sign-flip random seeds were fixed in the analysis plan. We report the direction and size of effects alongside adjusted inference; statistical significance alone is not interpreted as generalization across tasks. The original six contrasts remained the sole confirmatory analysis. V2 experiments were conducted independently as secondary robustness analyses.

For the frozen confirmation, secondary calibration diagnostics included mean absolute calibration error (MACE), the unweighted mean of absolute deviations between empirical and nominal interval coverage at 50\%, 80\%, 90\%, and 95\%. These descriptive checks used 1,280 fixed-state records; records are not independent inferential replicates. Gaussian negative log-likelihood means were not used for claims because their aggregate distribution was dominated by extreme values. The calibration exponent \(\alpha=0.5\) was selected on a development-stage fixed-state calibration diagnostic and frozen before confirmation.

\subsection{Secondary robustness analysis (V2)}
The V2 secondary analysis used a second randomized-tree-ensemble construction on an independently hashed D-drive snapshot of the same four benchmark tables and a separate namespace, protocol, split-membership file, seed series, and output directory; it is not a replication of the historical V5 confirmation. We used a bootstrapped \code{RandomForestRegressor} with 80 trees, minimum leaf size two, all features considered at each split, and single-threaded fitting. Raw uncertainty was the population standard deviation (\(ddof=0\)) of tree predictions. LOCAL uncertainty was derived from OOB residual ratios using an ExtraTrees local-scale model, five-fold cross-fitting, median normalization, clipping to [0.25,4], and the protocol-fixed geometric interpolation \(\alpha=0.5\). The four-task IID and chemical-system OOD splits used deterministic 20\% source, 20\% diagnostic, and 60\% target partitions with 20 fixed seeds. The 0.5\%, 1\%, and 5\% target-tail endpoints were post-processed from the same stored campaigns. Candidate-level target values were attached only after ranking and batch selection for retrospective evaluation; diagnostic labels were used only for descriptive calibration summaries. The campaign code also records MACE on the fixed diagnostic partition at nominal 50\%, 80\%, 90\%, and 95\% interval coverage for both RAW and LOCAL uncertainty at each stage. This auxiliary calibration diagnostic is summarized descriptively and is separate from the V2 utility endpoint.

\section{Results}
\subsection{Can a calibration transformation improve calibration without changing a rank decision?}
The independently executed rank-GLOBAL check yielded identical selected material identifiers to RAW in 16 of 16 task--split--seed checks. This is an implementation verification of the proposition, not a utility comparison; the check concerns rank acquisition.

\subsection{Does local rescaling perturb candidate selection?}
At the initial state of each campaign, LOCAL and RAW were evaluated on the same fitted surrogate and remaining candidate pool. The mean fraction of the RAW batch replaced under LOCAL was 0.255 for rank acquisition and 0.271 for magnitude acquisition; the corresponding medians were 0.248 and 0.252. These values summarize 320 task--split--seed initial states per acquisition rule. They are a post-analysis descriptive summary of stored decision diagnostics, not an additional confirmatory test or 320 independent inferential replicates. Supplementary Figure~S2 visualizes the means and medians. They establish that the intervention reached candidate selection at a shared state, while the campaign-level contrasts quantify the utility of its subsequent trajectory.

\subsection{Do changed selections improve average discovery utility?}
The mean seed-level \Lfourq was 0.5175 for rank-RAW and 0.5364 for rank-LOCAL. The paired LOCAL--RAW difference was 0.0190 (95\% bootstrap CI, 0.0122--0.0259; Holm-adjusted \(p_{\mathrm{adj}}=0.00030\)); 31 of 40 seed blocks (77.5\%) were positive. Under magnitude acquisition, mean \Lfourq was 0.6038 for RAW and 0.6182 for LOCAL. The paired difference was 0.0144 (95\% CI, 0.0083--0.0206; \(p_{\mathrm{adj}}=0.00030\)), with 29 of 40 positive seed blocks (72.5\%). Both point estimates exceeded the protocol-defined 0.005 threshold, and their reported bootstrap intervals also lay above that threshold. The sign-flip tests evaluated a zero mean difference, not a separate null hypothesis at 0.005.

The six prespecified LOCAL contrasts were positive and passed Holm-adjusted sign-flip tests (Table~\nef{tab:effects}). Supplementary Figure~S5 displays all 40 paired-seed differences for each contrast, showing the seed-level spread behind the pooled estimates. Comparisons with MATCHED support an advantage over the specific exploration-matching procedures used here; the comparison does not cover every possible acquisition policy.

\input{tables/primary_table1.tex}

\begin{figure}[ht]
\centering
\includegraphics[width=0.92\linewidth]{figures/primary_effects_origin_normalized.pdf}
\caption{Paired-seed estimates for the six prespecified primary contrasts in \Lfourq. Points show mean differences; horizontal bars show 95\% percentile-bootstrap confidence intervals. Blue and teal distinguish rank and magnitude acquisition, respectively. The left boundary of the horizontal scale is zero; positive values favor LOCAL. Estimates average equally over the eight fixed environments and do not imply a benefit in each environment.}
\label{fig:effects}
\end{figure}

\subsection{Is the utility effect consistent across properties?}
The positive pooled effects did not hold in every task. For \code{matbench\_mp\_gap}, mean LOCAL--RAW differences were $-$0.03823 under IID and $-$0.03465 under chemical OOD for rank acquisition; only 3/40 and 9/40 seed differences were positive. For magnitude acquisition, the corresponding means were $-$0.00195 and $-$0.00929, with 17/40 and 16/40 positive differences. The other three tasks had positive mean differences in both splits under both acquisition rules. The largest rank gains occurred in \code{matbench\_log\_kvrh} (log bulk modulus): 0.07365 under IID and 0.06382 under chemical OOD. Figure~\nef{fig:heterogeneity} visualizes these fixed-environment means; Supplementary Table~S2 reports the complete descriptive decomposition. No task-specific significance claim or causal explanation follows from these summaries.

\begin{figure}[tbp]
\centering
\includegraphics[width=\linewidth]{figures/Task_Split_Heterogeneity_normalized.pdf}
\caption{Task--split heterogeneity in LOCAL--RAW discovery utility. Blue squares and teal circles show mean paired differences in \Lfourq for rank and magnitude acquisition, respectively, over 40 seeds in each fixed environment. Symbols are displaced slightly vertically to separate the two policies. The gray line denotes zero; values to its left favor RAW. These are descriptive means with no error bars or environment-specific significance tests. All eight environments are shown, including both negative materials-project gap effects.}
\label{fig:heterogeneity}
\end{figure}

\subsection{Does improved initial-state MACE account for the MP-gap losses?}
Across 40 MP-gap paired seeds, LOCAL lowered stage-0 MACE by 0.0063 (IID) and 0.0088 (chemical OOD), yet mean LOCAL--RAW \Lfourq remained negative under both acquisition rules. Seed-block MACE--\Lfourq correlations were weak (rank, 0.03/0.08; magnitude, 0.14/0.02 for IID/OOD; Supplementary Table~S4 and Fig.~\nef{fig:mpgapdiagnostic}). In a separate diagnostic, MACE change correlated with first-batch disagreement in chemical OOD (Spearman 0.41 rank, 0.40 magnitude), but not IID (0.02, 0.00; Supplementary Fig.~S4). These exploratory associations do not evaluate candidate value or establish a mechanism.



\subsection{What happened to shared-state calibration?}
In the fixed-state diagnostic summary, OOF-LOCAL with \(\alpha=0.5\) had the lowest mean MACE (0.0412) among the reported conditions, compared with 0.0610 for RAW and 0.0465 for GLOBAL. Supplementary Figure~S1 compares the means and medians of the four principal calibration conditions, while Supplementary Figure~S3 shows mean observed coverage by nominal interval level. These are descriptive confirmation-state summaries, not independent hypothesis tests. They do not establish conditional calibration for every subgroup. The extreme-tailed Gaussian NLL mean is omitted from interpretation; its median and mean diverged sharply in the source summary, indicating sensitivity to a small number of extreme values.

\begin{figure}[H]
\centering
\includegraphics[width=0.94\linewidth]{figures/Exploratory_MP_Gap_MACE_Utility.pdf}
\caption{Exploratory diagnostic association between initial-state calibration and campaign utility for \code{matbench\_mp\_gap}. Each point is one paired seed: the horizontal coordinate is stage-0 shared-state MACE for OOF-LOCAL ($\alpha=0.5$) minus RAW, and the vertical coordinate is campaign \Lfourq for LOCAL minus RAW. Negative horizontal values indicate lower MACE; negative vertical values indicate lower utility. Diamonds mark the mean across 40 paired seeds, and the annotated Spearman coefficients are descriptive. The fixed zero lines show that rank acquisition has mean changes in the lower-left quadrant for both splits. These summaries neither identify displaced candidates nor support a chemical mechanism, and they are not additional hypothesis tests.}
\label{fig:mpgapdiagnostic}
\end{figure}

\subsection{Does a second randomized-tree-ensemble analysis preserve the primary direction?}
The second randomized-tree-ensemble analysis showed policy- and tail-dependent differences rather than a uniform LOCAL advantage (Fig.~\nef{fig:v2_tail_sensitivity}). After equal weighting across the eight task--split environments, the mean paired LOCAL--RAW difference at the top-1\% endpoint was +0.0057 for rank acquisition (13/20 seed blocks positive) and +0.0149 for magnitude acquisition (16/20 positive). At the top-0.5\% endpoint, the rank difference was near zero ($-$0.0002; 9/20 positive), whereas magnitude remained positive (+0.0183; 16/20 positive). At top 5\%, rank was positive (+0.0158; 17/20 positive) and magnitude negative ($-$0.0111; 5/20 positive). The error bars show the observed seed range, not confidence intervals; V2 was descriptive and no hypothesis tests were performed. Task--split means varied around these pooled summaries; the task--split heatmap and full descriptive table are provided in the Supplementary Information. Initial-batch disagreement averaged 25.8\% for rank and 34.6\% for magnitude, indicating that local rescaling perturbed decisions but did not guarantee a utility gain (see the V2 initial-batch disagreement table in the Supplementary Information). V2 is a secondary analysis using a second tree-ensemble construction, not a replication of the frozen primary study.

\begin{figure}[tbp]
\centering
\includegraphics[width=0.98\linewidth]{figures/v2_tail_robustness.pdf}
\caption{V2 sensitivity of LOCAL--RAW integrated-recall differences to the target-tail definition under (a) rank and (b) magnitude acquisition. Points are the equal-weight means across eight task--split environments, averaged over 20 paired seed blocks. Error bars span the observed minimum to maximum of those 20 seed-level means; they are not confidence intervals. Positive values favor LOCAL. The original top-1\% endpoint remains the primary endpoint; the 0.5\% and 5\% results are secondary descriptive sensitivity analyses, with no V2 hypothesis tests.}
\label{fig:v2_tail_sensitivity}
\end{figure}

\subsection{Do changed V2 candidates enter the intended target tail?}
Across the eight fixed environments, the equal-weight LOCAL-only minus RAW-only top-1\% tail-fraction difference was +0.0007 for rank acquisition and +0.0002 for magnitude acquisition (20 paired seeds and four stages per environment). The MP-gap chemical-OOD rank exchange had lower top-1\% enrichment under LOCAL (0.0265 versus 0.0305; difference $-0.0040$, ratio 0.87), whereas the IID rank difference was near zero and the magnitude differences were small in both splits. In contrast, rank acquisition on log-GVRH (log shear modulus) chemical OOD showed a positive enrichment difference (+0.0053). These retrospective summaries are consistent with poorer tail alignment in one MP-gap policy--split condition, but not a common explanation for the MP-gap utility losses. Mean target value and tail enrichment answer different questions; no causal interpretation or V2 significance claim is made (see the candidate-tail-enrichment table and figure in the Supplementary Information).

\subsection{Is the decision effect sensitive to the strength of local rescaling?}
Across four prespecified environments and 20 paired V2 seeds, the equal-weight top-1\% LOCAL--RAW means at $\alpha=0.25$, 0.50, 0.75, and 1.00 were $-$0.0031, $-$0.0121, $-$0.0181, and $-$0.0163 for rank, and +0.0026, +0.0121, $-$0.0121, and $-$0.0450 for magnitude. At 5\%, rank stayed positive (+0.0138 to +0.0222), while magnitude declined from $-$0.0008 to $-$0.0605. Seed ranges are descriptive, with no confidence intervals or tests. The frozen primary remains $\alpha=0.5$; this assesses robustness, not tuning (Supplementary Table~S9; Fig.~S7).

\section{Discussion}
\subsection{Calibration metrics do not determine decisions}
The results separate three outcomes that are often discussed together. The shared-state diagnostics describe interval calibration; the initial batch comparison shows whether local rescaling reaches the acquisition decision; and the paired campaigns measure target-tail utility. Each answers a distinct question. In particular, a lower MACE cannot imply a different ranking, and a changed ranking cannot establish a better query sequence.

\subsection{Acquisition representation determines sensitivity}
The rank proposition gives a policy-specific invariance: a positive scalar changes the magnitude of uncertainty but preserves its within-pool ordering. The 16 trajectory matches verify that implementation property. Candidate-dependent rescaling can alter relative rankings, while magnitude acquisition also responds to the numerical uncertainty scale. Decision relevance therefore depends on the interaction between the transformation and the acquisition representation.

\subsection{Perturbation did not guarantee utility}
Across the eight fixed environments, local candidate-dependent rescaling improved average utility under both acquisition rules and exceeded the prespecified PERMUTED and MATCHED controls. The result is conditional on this benchmark protocol. Both \code{matbench\_mp\_gap} splits showed lower mean utility under LOCAL than RAW. A retrospective, ID-matched audit now shows that the mean target value of same-stage LOCAL entrants was higher than that of RAW-displaced candidates under rank acquisition (+0.0673 IID; +0.1249 chemical OOD) but slightly lower under magnitude acquisition (-0.0162; -0.0143; Supplementary Table~S7). These mixed label shifts do not account for the lower top-1\% utility under both rules: mean property value among exchanged candidates is not the same as enrichment in the target-pool top tail. We therefore make no mechanistic or chemical claim.

The exploratory initial-state MACE--utility correlations were weak across MP-gap split-policy combinations (\(\nho=0.02\)--0.14; Fig.~\nef{fig:mpgapdiagnostic}) and do not explain the later campaign decline. The initial MACE--first-batch-disagreement correlation was moderate only in chemical OOD (\(\nho=0.40\)--0.41) and near zero in IID; it describes co-variation between two initial-state summaries, not whether changed candidates helped. The historical confirmation archive supports an ID-based mean-target-value displacement analysis but does not retain candidate prediction vectors or the frozen target-pool membership needed to measure its entrant-tail enrichment. The separate V2 candidate-state files do retain these fields and allow the exploratory tail analysis above. This distinction removes the data limitation for V2 without supplying a causal explanation of the primary MP-gap result.

The second randomized-tree-ensemble analysis adds an important boundary to the pooled primary result. Although the top-1\% mean differences remained positive under both acquisition rules, the direction changed with the tail definition: at 5\%, magnitude acquisition favored RAW on average while rank acquisition favored LOCAL. The task--split decomposition also retained substantial heterogeneity. The sequential alpha sensitivity further shows that the sign of top-1\% utility varies with intervention strength, with more negative magnitude effects at larger alpha. These descriptive secondary results do not overturn the frozen primary endpoint, and they do not support a general claim that local rescaling improves discovery for every property objective or tail definition.

The positive pooled effects are consistent with the direction of the materials-optimization improvements reported by Jiang et al. under globally temperature-scaled bootstrap uncertainty and expected improvement \citep{jiang2026calibration_optimization}. Harish reports a complementary boundary in rare-earth extraction: conformal calibration improved uncertainty reliability but left rank-based acquisition unchanged, while discovery rankings depended on the target definition \citep{harish2026objective}. The negative MP-gap effects show that this direction does not transfer uniformly across the task--split environments tested here. This is not a head-to-head comparison: the studies use different materials datasets, uncertainty estimators, acquisition policies, and utility definitions. Our results add evidence about policy sensitivity and task-level variation rather than a general ranking of calibration procedures.

\subsection{Implications and scope}
For scientific screening, evaluate a calibration intervention in sequence: determine whether the acquisition rule is sensitive to it, measure how candidate selection changes, and assess downstream utility for the intended property objective. The confidence intervals quantify algorithmic variability over paired seeds for these fixed benchmark environments; they should not be interpreted as uncertainty over a population of materials-property tasks. This study covers four inherited Matbench tasks, a composition-based representation, randomized-tree ensembles, two prescribed splits, and rank and magnitude acquisition in an offline simulation. No separate pre-confirmation rationale for task selection is recoverable, so task-level patterns remain descriptive and the results do not establish generalization beyond the tested fixed environments. The chemical-system split represents one defined form of shift; no physical synthesis or laboratory feedback loop was tested.

\section{Conclusion}
Calibration quality, acquisition sensitivity, candidate perturbation, and task-specific utility are distinct outcomes. Positive global scaling is decision-invariant under the specified rank rule. Candidate-dependent rescaling improved the frozen primary pooled endpoint but reduced utility on both MP-gap splits; in the second randomized-tree-ensemble analysis, utility also varied with policy, target-tail definition, and intervention strength. Retrospective entrant enrichment was poorer for rank acquisition on MP-gap chemical OOD but nearly unchanged for the other MP-gap policy--split combinations, so it does not explain the losses as a common mechanism. Calibration quality alone does not determine the scientific decision value of an uncertainty intervention.

\section*{Funding}
This research received no specific grant from any funding agency in the public, commercial, or not-for-profit sectors.

\section*{Data and code availability}
Analysis code, configuration files, aggregate result tables, verification manifests, and reproducibility documentation are available at \url{https://github.com/9liforeverloves-lgtm/decision-relevance-uncertainty-calibration}. The source Matbench datasets are available from their original providers and are not redistributed by this repository where redistribution rights are not established. The repository documents the data-acquisition and analysis workflow used for the reported results.

\fontsize{8}{9}\selectfont\setlength{\bibsep}{0pt}
\bibliographystyle{plainnat}
\bibliography{bibliography}
\end{document}
